%% GENERATED by python/paper/build.py from results/ -- do not edit by hand.
%% Rebuild:  .venv\Scripts\python.exe python\paper\build.py --only OECD26_Slopes_cses
%% Source:   results/oecd26/layer1.csv, layer2.csv, anchor.csv; data/oecdFigure1.csv, oecd26_theta.csv
\begin{table}[!htb]
\centering
\begin{threeparttable}
\caption{Figure \ref{fig:US:OECDdata}\textquotesingle s relations in the data and in the model, without Luxembourg}
\label{table:OECD26:slopes_cses}
\renewcommand{\arraystretch}{1.25}
\begin{tabularx}{\textwidth}{lYYY}
\toprule
\textbf{Relation} & \textbf{Data} & \textbf{Model} & \textbf{Own $\omega$} \\
\midrule 
Pension spending on population growth & $-14.91$\newline$(4.92)$ & $-20.35$\newline$(2.51)$ & $-15.82$\newline$(4.81)$ \\ % row: spending_growth
Pension spending on the Gini & $14.79$\newline$(26.28)$ & $-7.28$\newline$(23.70)$ & $14.10$\newline$(26.48)$ \\ % row: spending_gini
Pension spending on the Bismarckian index & $16.90$\newline$(3.94)$ & $65.25$\newline$(6.03)$ & $56.32$\newline$(6.51)$ \\ % row: spending_index
Bismarckian index on population growth & $-0.77$\newline$(0.18)$ & $-0.29$\newline$(0.03)$ & $-0.30$\newline$(0.07)$ \\ % row: index_growth
Bismarckian index on the Gini & $0.80$\newline$(1.10)$ & $0.12$\newline$(0.34)$ & $0.35$\newline$(0.42)$ \\ % row: index_gini
\bottomrule
\end{tabularx}
\begin{tablenotes}[flushleft]
\footnotesize
\item[] \textit{Note:} Least-squares slopes, standard errors in parentheses, 20 countries. The model\textquotesingle s index is $1/(2-\theta)$ of the design in force in 2020, chosen in 1990, and its spending $\tau(1-\alpha)$ in 2020 on the path of the design it chooses. The model: every country at the U.S.\ $\beta = 0.65$ and $\varepsilon = 13.09$, at $\omega = 1.20$, the mean of the countries\textquotesingle{} own $\omega$, each at its cost $\lambda_{n,t} = \varepsilon\,\Pi_{n,t}/R_{n,t}$; own $\omega$: each country\textquotesingle s $\omega$ calibrated to its spending at its observed design. Turnout: each country\textquotesingle s own propensities to vote by income group, from the CSES.
\end{tablenotes}
\end{threeparttable}
\end{table}
